\documentclass[11pt,a4paper]{scrartcl}

\usepackage{../manostilius_lt}

\title{Daugianaris}
\author{Andrius Berniukevičius}

\begin{document}
\maketitle


\section{Algebrinė suma}
\begin{definition}
\vocab{Algebrinė suma} -- tai kelių narių suma. Jos nariai sujungiami sudėties ženklu.	
\end{definition} 
Jeigu reiškinyje matome minuso ženklą, jo nelaikome atimties veiksmu. Minusas parodo, kad pridedamas neigiamas narys. Pavyzdžiui, reiškinį
\[
3a^2-5a+7
\]
pertvarkome taip:
\[
3a^2-5a+7=3a^2+(-5a)+7.
\]
Taigi tai yra algebrinė suma, kurios nariai yra $3a^2$, $-5a$ ir $7$. Joje sudedame tris narius, o antrasis iš jų yra neigiamas.

\section{Daugianario apibrėžimas}

Ankstesnėje temoje nagrinėjome vienanarius. Daugianariai sudaromi iš vienanarių, juos sudedant. Neigiamas narys gali būti užrašytas su minuso ženklu.

\begin{definition}
\vocab{Daugianaris} -- tai algebrinė vienanarių suma.
\end{definition}

Pavyzdžiui, šie reiškiniai yra daugianariai:
\[
3x^2-5x+7, \qquad a^2b-4ab^2+6, \qquad 8, \qquad -2xy^3.
\]

Vienanaris taip pat laikomas daugianariu, nes jį sudaro vienas vienanaris.

\begin{example}
Nustatykime, kurie reiškiniai yra daugianariai:
\begin{tasks}[style=itemize](2)
	\task $5x^3-2x+1$
	\task $\frac{3}{x}+4$
	\task $7a^2b-5ab^2$
	\task $\sqrt{x}+2$
	\task $-4m^2n$
	\task $3p^{-2}+p$
\end{tasks}
\textbf{Sprendimas}\\
Daugianariai yra $5x^3-2x+1$, $7a^2b-5ab^2$ ir $-4m^2n$. Jie yra sudaryti iš vienanarių.

Reiškinys $\frac{3}{x}+4$ nėra daugianaris, nes kintamasis yra vardiklyje. Reiškinys $\sqrt{x}+2$ nėra daugianaris, nes kintamasis yra po šaknies ženklu. Reiškinys $3p^{-2}+p$ nėra daugianaris, nes kintamojo $p$ rodiklis yra neigiamas.
\end{example}

\section{Daugianario nariai}

\begin{definition}
\vocab{Daugianario nariai} -- tai vienanariai, iš kurių sudarytas daugianaris. Juos skiria sudėties arba atimties ženklai.
\end{definition}

\begin{example}
Panagrinėkime daugianarį:
\[
4a^2b-3ab^2+7.
\]
\textbf{Sprendimas}\\
Šio daugianario nariai yra $4a^2b$, $-3ab^2$ ir $7$.
\end{example}

\begin{remark}
Atimties ženklą patogu suprasti kaip neigiamo vienanario pridėjimą. Pavyzdžiui:
\[
x^2-5x+3=x^2+(-5x)+3.
\]
Taigi daugianario nariai yra $x^2$, $-5x$ ir $3$.
\end{remark}

\section{Dvinaris ir trinaris}

\begin{definition}
\vocab{Dvinaris} -- tai daugianaris, sudarytas iš dviejų vienanarių.
\end{definition}

Pavyzdžiui, $3x^2-5x$ yra dvinaris. Jo nariai yra $3x^2$ ir $-5x$.

\begin{definition}
\vocab{Trinaris} -- tai daugianaris, sudarytas iš trijų vienanarių.
\end{definition}

Pavyzdžiui, $2a^2b-3ab^2+7$ yra trinaris. Jo nariai yra $2a^2b$, $-3ab^2$ ir $7$.

\begin{example}
Nustatykime, kurie reiškiniai yra dvinariai, o kurie -- trinariai:
\begin{tasks}[style=itemize](2)
	\task $x^2+4$
	\task $5a^2-3a+1$
	\task $-2mn$
	\task $p^3-p^2+p$
\end{tasks}
\textbf{Sprendimas}\\
Reiškinys $x^2+4$ yra dvinaris, nes turi du narius. Reiškiniai $5a^2-3a+1$ ir $p^3-p^2+p$ yra trinariai, nes kiekvienas turi po tris narius. Reiškinys $-2mn$ yra vienanaris, todėl jis nėra nei dvinaris, nei trinaris.
\end{example}

\section{Kada reiškinys nėra daugianaris}

Kad reiškinys būtų daugianaris, kiekvienas jo narys privalo būti vienanaris. Todėl reiškinys nėra daugianaris, jeigu jame yra nors vienas iš šių atvejų:
\begin{itemize}
	\item kintamasis vardiklyje, pavyzdžiui, $\frac{5}{x}$;
	\item kintamasis su neigiamu laipsnio rodikliu, pavyzdžiui, $2x^{-3}$;
	\item kintamasis po šaknies ženklu, pavyzdžiui, $\sqrt{y}$;
	\item kintamasis su trupmeniniu laipsnio rodikliu, pavyzdžiui, $a^{\frac{1}{2}}$.
\end{itemize}

\begin{example}
Nustatykime, ar reiškinys yra daugianaris:
\[
(x+3)(x-2).
\]
\textbf{Sprendimas}\\
Atskliaudžiame:
\[
(x+3)(x-2)=x^2-2x+3x-6=x^2+x-6.
\]
Gautas reiškinys yra vienanarių suma, todėl $(x+3)(x-2)$ yra daugianaris.
\end{example}

\begin{example}
Nustatykime, ar reiškinys yra daugianaris:
\[
\frac{x^2-1}{x+1}.
\]
\textbf{Sprendimas}\\
Nors trupmeną galima suprastinti, pradiniame reiškinyje vardiklyje yra kintamasis. Todėl šis reiškinys nėra daugianaris.
\end{example}

\section{Uždaviniai}

\begin{problem}
Nustatykite, kurie reiškiniai yra daugianariai.
\begin{tasks}[label={(\arabic*)}, label-offset=0.9em](2)
	\task $6x^2-3x+8$
	\task $\frac{4}{y}-2$
	\task $-5a^3b+2ab^2-7$
	\task $\sqrt{m}+m$
	\task $9$
	\task $x^{-2}+3x$
	\task $\frac{3}{5}p^2q-4q$
	\task $\frac{a+2}{7}$
	\task $2x\sqrt{y}-1$
	\task $(a-1)(a+4)$
	\task $\frac{5m^2}{n}+3$
	\task $-7r^4s^2$
	\task $b^{\frac{2}{3}}+b$
	\task $(x+2)^2$
\end{tasks}
\end{problem}

\begin{problem}
Išvardykite daugianario narius.
\begin{tasks}[label={(\arabic*)}, label-offset=0.9em](2)
	\task $3x^2-7x+4$
	\task $-5a^3b+2ab^2-9$
	\task $m^2n-6mn^2+n$
	\task $8p^4$
	\task $-x^3+4x^2-x+1$
	\task $\frac{2}{3}a^2b-5b+6$
\end{tasks}
\end{problem}

\newpage
\answerssection

\begin{problem} \phantom{text}
\begin{tasks}[label={(\arabic*)}, label-offset=0.9em](2)
	\task daugianaris;
	\task ne daugianaris;
	\task daugianaris;
	\task ne daugianaris;
	\task daugianaris;
	\task ne daugianaris;
	\task daugianaris;
	\task daugianaris;
	\task ne daugianaris;
	\task daugianaris;
	\task ne daugianaris;
	\task daugianaris;
	\task ne daugianaris;
	\task daugianaris.
\end{tasks}
\end{problem}

\begin{problem} \phantom{text}
\begin{tasks}[label={(\arabic*)}, label-offset=0.9em](2)
	\task $3x^2$, $-7x$, $4$;
	\task $-5a^3b$, $2ab^2$, $-9$;
	\task $m^2n$, $-6mn^2$, $n$;
	\task $8p^4$;
	\task $-x^3$, $4x^2$, $-x$, $1$;
	\task $\frac{2}{3}a^2b$, $-5b$, $6$.
\end{tasks}
\end{problem}

\end{document}
